\documentclass[a4paper]{article}

\usepackage{graphicx}
\usepackage{amsmath,amssymb}
\usepackage{minted}
\usepackage{multirow}
\usepackage{float}
\usepackage{algorithm}
\usepackage{algpseudocode}
\usepackage{todonotes}
\usepackage{forest}
\usepackage{xstring}
\usepackage{xcolor}
\usepackage[most]{tcolorbox}
\usepackage{xparse}
\usepackage{natbib}

\usetikzlibrary{graphs,graphdrawing}
\usegdlibrary{layered}

\title{PDF Preview Page Boundary Scrolling}
\date{2026-09-08}
\author{Mohammad Rahmani}

\begin{document}
   \maketitle

   PDF Preview scrolling, panning, zoom anchoring, and explicit scroll bars now use the rendered PDF page rectangle as their movement boundary instead of the larger internal \texttt{PdfMultiPageView} content extent. This prevents the viewport from stopping before the true left, right, top, or bottom edge of a tight-page PDF and prevents movement beyond the actual page geometry.

   The boundary calculation maps all four corners of the current rendered page into the scrollable content coordinate system and derives one shared horizontal and vertical range for manual panning, scroll-bar positioning, view restoration, and zoom-anchor restoration.

   \bibliographystyle{plainnat}
   \bibliography{/home/donkarlo/Dropbox/repo/nd_shared_research/bibliography/bibliography.bib}
\end{document}
